\documentclass[11pt]{article}

\usepackage[T1]{fontenc}
\usepackage[utf8]{inputenc}
\usepackage[margin=1in]{geometry}
\usepackage{amsmath,amssymb,amsthm,mathtools}
\usepackage{array}
\usepackage{enumitem}
\usepackage{listings}
\usepackage{microtype}
\usepackage[hidelinks]{hyperref}

\newtheorem{theorem}{Theorem}
\newtheorem{lemma}{Lemma}
\newtheorem{proposition}{Proposition}
\newtheorem{definition}{Definition}
\newtheorem{claim}{Claim}
\newtheorem{remark}{Remark}
\newtheorem{corollary}{Corollary}

\newcommand{\NS}{\mathrm{NS}}
\newcommand{\CM}{\mathrm{CM}}
\newcommand{\Pack}{\mathrm{Pack}}
\newcommand{\Part}{\mathrm{Part}}
\newcommand{\Field}{\mathrm{Field}}
\newcommand{\Member}{\mathrm{Member}}
\newcommand{\Exit}{\mathrm{Exit}}
\newcommand{\Legal}{\mathrm{Legal}}
\newcommand{\Pass}{\mathrm{Pass}}
\newcommand{\Fail}{\mathrm{Fail}}
\newcommand{\Owork}{\mathfrak O_{\NS}^{\mathrm{work}}}
\newcommand{\WitnessRecord}{\mathcal W}
\newcommand{\Rterm}{\mathcal R_*}
\newcommand{\Ffail}{\mathcal F}
\newcommand{\Tthree}{\mathbb T^3}
\newcommand{\Rthree}{\mathbb R^3}
\newcommand{\Hs}{H^s}
\newcolumntype{L}[1]{>{\raggedright\arraybackslash}p{#1}}

\setlength{\emergencystretch}{1em}

\title{Three-Dimensional Navier--Stokes:\\
       Same-Solution Contrapositive and Whole-Space Class Exit}
\author{Thomas Birnie}
\date{June 7, 2026}

\begin{document}
\maketitle

\begin{abstract}
This manuscript develops the same-solution contrapositive route for the three-dimensional Navier--Stokes smoothness problem.
The core object is not a new positive estimate for every derivative.
It is a test of any alleged finite terminal obstruction for the original smooth solution.
The obstruction has to remain attached to a positive same-fluid packet, the same Navier--Stokes equation has to keep participating on that packet, and one positive-scale field readout has to stay coherent.
The ordered faces of that test are
\[
  \Pack_Q,\qquad \Part_{N,Q},\qquad \Field_{N,r,Q}.
\]
The first failed face supports the class-exit conclusion
\[
  \Exit(Q;\Owork):=\neg\Member(Q;\Owork).
\]
The periodic \(\Tthree\) route is carried by the installed same-solution witness chain: terminal entry, Pack/Part/Field face exhaustion, and class-exit embedding.
The whole-space \(\Rthree\) branch is read by the same test after the initial tail, compact-core far-kernel source, and bounded exterior frequencies have been separated.
Any surviving exterior high-frequency source is dyadic scale escape, and the ordered test lands that survivor in Pack, Part, or Field.
\end{abstract}

\section{The Problem}

Let \(M\) be either \(\Tthree\) or \(\Rthree\).
On \(\Rthree\), the datum is smooth, divergence-free, and rapidly decaying.
Let \(u:M\times[0,T)\to\mathbb R^3\) and \(p:M\times[0,T)\to\mathbb R\) satisfy
\[
  \partial_t u+(u\cdot\nabla)u+\nabla p=\nu\Delta u,
  \qquad
  \nabla\cdot u=0,
\]
with fixed \(\nu>0\), zero external force, and smooth divergence-free initial datum
\[
  u(\cdot,0)=u_0.
\]
The smoothness question asks whether the smooth solution determined by \(u_0\) can stop at a finite time because smoothness breaks down.

The ordinary direct route tries to prove estimates strong enough to keep every derivative finite.
The class-membership route begins from the opposite end.
It takes a claimed finite terminal obstruction and asks whether that obstruction is still a legal terminal object for the same Navier--Stokes evolution.

\section{Why The Direct Estimate Route Stalls}

Energy gives a real law.
For smooth divergence-free periodic solutions,
\[
  \frac12\frac{d}{dt}\|u(t)\|_{L^2}^2
  +\nu\|\nabla u(t)\|_{L^2}^2=0.
\]
This law controls kinetic energy and the time integral of the gradient.
It does not by itself control all derivatives.

One derivative higher, the enstrophy calculation exposes the familiar obstruction:
\[
  \frac12\frac{d}{dt}\|\nabla u(t)\|_{L^2}^2
  +\nu\|\Delta u(t)\|_{L^2}^2
  =
  \int_{\Tthree}(u\cdot\nabla)u\cdot\Delta u\,dx.
\]
The left side contains viscous dissipation.
The right side contains the nonlinear source.
A direct global proof would need a closed estimate that absorbs the dangerous part of that source without asking for exactly the higher control the argument is trying to prove.

This is why the contrapositive route is used here.
The paper does not pretend that the enstrophy estimate has been closed.
It asks what a claimed terminal breakdown would have to be as an object of the original equation.

\section{Contrapositive Reorientation}

Assume a finite terminal obstruction is claimed for the original smooth datum.
The claim has to point to something at the alleged terminal time.
That something is the terminal witness under test.

The class-membership question is narrower than a new positive estimate:
\begin{quote}
Does the alleged terminal witness remain a legal same-fluid Navier--Stokes object on a positive material packet, with the same equation participating, and with one coherent positive-scale field readout?
\end{quote}

The pass branch is the branch that remains a member of the working class.
The fail branch is not a new nonsmooth solution inside the class.
It is the first place where the alleged terminal object loses a required same-solution face.
The proof therefore uses the split
\[
  \Pass\Longrightarrow \Member(Q;\Owork),
  \qquad
  \Fail\Longrightarrow \Exit(Q;\Owork).
\]

\section{The Working Class}

\begin{definition}[Working same-fluid class]
The working class \(\Owork\) is the paper's class-membership object for the periodic Navier--Stokes evolution from \(u_0\).
A terminal record belongs to the class only when it remains attached to the same evolution and retains the required packet, participation, and field faces.
\end{definition}

\begin{definition}[Pack face]
\(\Pack_Q\) means that the terminal object has a positive same-fluid material packet for the selected record \(Q\).
Failure of \(\Pack_Q\) means the alleged terminal object has lost the material body needed to be tested as the same solution.
\end{definition}

\begin{definition}[Participation face]
\(\Part_{N,Q}\) means that, through depth \(N\), the same Navier--Stokes equation acts on the selected packet.
This includes the same pressure-viscosity participation, not a neighboring equation, filtered shadow, or exported surrogate.
\end{definition}

\begin{definition}[Field face]
\(\Field_{N,r,Q}\) means that, after Pack and Part are fixed, one positive scale \(r>0\) still carries coherent field data through depth \(N\).
Failure of every positive scale is a field-side class failure.
\end{definition}

\begin{definition}[Class exit]
The class-exit conclusion is
\[
  \Exit(Q;\Owork):=\neg\Member(Q;\Owork).
\]
It is supported by a face failure.
It is not identical to a face failure.
\end{definition}

\section{Terminal Witness Consumption}

The audited route lock states the branch law in this order:
\[
\begin{aligned}
&\text{finite Clay terminal witness}\\
&\Longrightarrow \text{CM-test entry}\\
&\Longrightarrow
  \neg\Pack_Q\ \vee\ \neg\Part_{N,Q}\ \vee\
  \forall r>0\,\neg\Field_{N,r,Q}\\
&\Longrightarrow \Exit(Q;\Owork).
\end{aligned}
\]

The order matters.
Pack is tested first because the witness needs a positive same-fluid carrier.
Part is tested after a carrier has survived because the same Navier--Stokes equation must act on that carrier.
Field is tested after Pack and Part because positive-scale coherence is meaningful only on a participating same-fluid packet.

\begin{lemma}[First failed face gives class exit]
Let \(Q\) be a terminal record admitted to the CM test.
When the first failed requirement in the ordered Pack/Part/Field test is
\[
  \neg\Pack_Q,
  \qquad
  \Pack_Q\wedge\neg\Part_{N,Q},
  \qquad\text{or}\qquad
  \Pack_Q\wedge\Part_{N,Q}\wedge\forall r>0\,\neg\Field_{N,r,Q},
\]
the record supports \(\Exit(Q;\Owork)\).
\end{lemma}

\begin{proof}
Membership in \(\Owork\) requires the terminal record to keep the same-fluid packet, same-equation participation, and positive-scale field readout.
The first failed face removes one required condition for membership while preserving the ordered information needed to say where the failure occurred.
By definition of class exit, that loss supports \(\neg\Member(Q;\Owork)\), written as \(\Exit(Q;\Owork)\).
\end{proof}

\section{No Third Branch}

The pass branch and fail branch are not two independent solution theories.
They are the two outcomes of testing the same alleged terminal object.
The pass branch remains a legal member branch.
The fail branch carries a face failure and therefore exits the working class.

\begin{proposition}[No in-class nonsmooth third branch]
A terminal record cannot simultaneously be the same legal \(\Member(Q;\Owork)\) continuation branch and the same record carrying the derived terminal Pack/Part/Field face failure.
\end{proposition}

\begin{proof}
The member branch retains the faces required by \(\Owork\).
The failure branch loses at least one of those faces in the ordered terminal test.
The same record cannot retain and lose the same required membership structure at once.
\end{proof}

\section{The Periodic CM Chain}

The periodic proof does not rest on a separate positive estimate for every branch family.
It rests on the same-solution terminal object.
The transported whole-torus atlas \(Q_{\mathrm{atlas}}\) admits the alleged finite terminal witness to the CM test and asks only three questions:
\[
  \Pack_Q,\qquad \Part_{N,Q},\qquad \Field_{N,r,Q}.
\]
If all three survive, the same branch has the \(H^s\), \(s>5/2\), continuation readout.
If one fails first, the failed face is the class-exit witness
\[
  \Exit(Q;\Owork):=\neg\Member(Q;\Owork).
\]
The installed terminal-witness theorem says there is no fourth same-solution terminal outcome outside those three faces.

\section{Periodic Completion}

\begin{theorem}[Periodic CM completion]
For the periodic \(\Tthree\) Navier--Stokes statement, a finite same-solution terminal breakdown cannot remain as an in-class nonsmooth third branch.
\end{theorem}

\begin{proof}
Let \(Q\) be the alleged finite terminal record.
By terminal entry, \(Q\) is the same-fluid object under the CM test.
The ordered test has two possible outcomes.
On the pass side, \(Q\) retains Pack, Part, and a positive Field scale, so the \(H^s\) continuation criterion extends the same solution.
On the fail side, the first failed face is Pack, Part, or Field, and the CM embedding reads that face as \(\Exit(Q;\Owork)\).
Thus the alleged terminal record is either the same smooth continuation branch or a classified class-exit witness.
It is not a legal in-class nonsmooth terminal object.
\end{proof}

\section{Support Families}

Older positive-forward material remains useful only after it is given the right proof role.
Source-reserve estimates, selector arguments, terminal tail routes, pressure-source diagnostics, Euler comparisons, and whole-space export attempts may explain how a terminal record was found or pressure-tested.
They become proof-bearing in this draft only when they land a same-record statement in \(\Pack_Q\), \(\Part_{N,Q}\), or \(\Field_{N,r,Q}\), or when they break terminal entry, finite-failure exhaustion, or CM embedding.

This quarantine rule prevents a common overclaim.
A support theorem can look powerful while proving a statement about a surrogate, a filtered object, a downstream export, or a different route.
The proof uses support only after the support has been translated onto the same terminal witness.

\section{Whole-Space Class Exit}

The Clay problem is posed on \(\Rthree\).
The periodic route therefore has to be exported without changing the theorem.
The current export problem is not a vague request for ``more decay''.
It has a specific Duhamel form.

Let \(u=e^{t\nu\Delta}u_0+B(u,u)\), with
\[
  B(u,u)(t)
  =
  -\int_0^t e^{(t-s)\nu\Delta}
  \mathbb P\nabla\cdot(u\otimes u)(s)\,ds.
\]
Outside a large ball, the rapid-decay part of the initial heat flow is harmless.
The compact-core portion of the nonlinear source also has far-kernel decay.
The remaining issue is the genuinely exterior nonlinear source.

\begin{theorem}[Whole-space export reduction]
Assume the periodic same-solution Pack/Part/Field classification is available for the local terminal packet.
After the initial heat tail and compact-core nonlinear source are removed, the genuinely exterior nonlinear Duhamel source has the following split.
If its exterior \(H^s\) tail vanishes, the whole-space terminal object has the same class-membership split as the periodic object at the selected packet.
If the exterior \(H^s\) source survives, then bounded frequencies vanish by exterior \(L^2\) tightness, the survivor escapes to dyadic scale, and the branch lands in Pack, Part, or Field.
\end{theorem}

\begin{proof}
The terminal test is local on the selected positive packet once Pack has supplied the carrier and Part has supplied same-equation participation.
The initial heat-flow tail cannot create a new terminal face at infinity because rapid decay is preserved by the heat kernel.
The compact-core nonlinear source has far-field smoothing through the heat kernel and therefore cannot create the missing exterior \(H^s\) obstruction.
After those two pieces are removed, the only possible difference between the periodic readout and the \(\Rthree\) readout is the exterior nonlinear Duhamel source.
The stated tail bound removes that difference on the direct export branch.
On the survivor branch, exterior \(L^2\) tightness removes bounded frequencies.
The remaining dyadic survivor has no positive material address, no same-fluid participation, or, after Pack and Part are retained, no fixed positive Field scale.
The same Pack/Part/Field face test therefore applies to the whole-space terminal object.
\end{proof}

\begin{lemma}[Exterior dyadic survivor face landing]
Let \(W_R\) be the exterior nonlinear Duhamel contribution after the initial heat tail and fixed compact-core source have been removed.
Assume a terminal exterior survivor remains: for some \(s>5/2\), some \(t_j\uparrow T_*\), some \(R_j\to\infty\), and dyadic \(N_j\to\infty\),
\[
  N_j^s\|P_{N_j}W_{R_j}(t_j)\|_{L^2}
\]
stays bounded below.
Then the survivor supplies one of the CM face failures
\[
  \neg\Pack_Q,\qquad
  \neg\Part_{N,Q},\qquad
  \forall r>0\,\neg\Field_{N,r,Q}.
\]
\end{lemma}

\begin{proof}
Exterior \(L^2\) tightness removes the bounded-frequency part of an exterior tail.
The survivor is therefore a genuine dyadic escape.

First test Pack.
A terminal exterior source with no positive same-fluid material address has no packet on which the original solution can be tested, so it is \(\neg\Pack_Q\).

Assume Pack survives.
Then test Part.
If the packet is not carried by the same Navier--Stokes pressure, viscosity, and source ancestry, the survivor is \(\neg\Part_{N,Q}\).

Assume Pack and Part survive.
Then the survivor is on one same participating packet.
A positive Field scale would give a fixed coherent field window and a finite derivative tower above the \(H^s\) continuation level.
On that fixed witness record, Littlewood--Paley tails vanish at high dyadic frequency:
\[
  \lim_{N\to\infty} N^s\|P_N W_R(t)\|_{L^2}=0.
\]
That contradicts the lower bound along \(N_j\to\infty\).
Hence no fixed positive Field scale survives, so \(\forall r>0\,\neg\Field_{N,r,Q}\).
\end{proof}

\begin{theorem}[Whole-space CM completion]
For smooth rapidly decaying divergence-free data on \(\Rthree\), a finite terminal whole-space branch has the same CM pass-or-exit conclusion as the periodic branch.
The branch either keeps the exterior \(H^s\), \(s>5/2\), continuation tail, or the surviving exterior source is consumed by Pack, Part, or Field and embedded as
\[
  \Exit(Q;\Owork):=\neg\Member(Q;\Owork).
\]
\end{theorem}

\begin{proof}
The Duhamel split leaves only the genuinely exterior nonlinear source after the harmless heat tail and fixed compact-core source have been removed.
When the exterior \(H^s\) tail vanishes, the whole-space branch reads like the periodic branch on the selected terminal packet and continues by the standard \(H^s\) criterion.
When it survives, the preceding lemma supplies the first failed face.
The CM embedding reads that face failure as class exit.
Thus a finite terminal whole-space claim is either the smooth member branch or a class-exit witness.
It does not produce an in-class nonsmooth third branch.
\end{proof}

\begin{remark}
The direct exterior nonlinear source estimate is a stronger positive route.
The proof above does not use that positive tail estimate as an assumption.
When tail vanishing is not supplied, the survivor branch is still tested as the same alleged terminal object and is classified by Pack, Part, or Field.
A separate positive proof of tail vanishing would have to come from the equation and the terminal packet structure, not from global smoothness, the full Beale--Kato--Majda criterion, or a global bound equivalent to the desired regularity conclusion.
\end{remark}

\section{Terminal Records Are Same-Solution Objects}

The proof uses the word terminal in a narrow way.
It does not mean that a bad-looking distribution has been found near the alleged terminal time.
It means that the alleged obstruction is still being offered as evidence about the original solution.
That is why the same-solution requirements are part of the theorem rather than explanatory decoration.

\begin{definition}[Selected terminal record]
A selected terminal record \(Q\) is the data extracted from the original smooth solution along a sequence \(t_j\uparrow T_*\), together with the packet and scale information needed to test whether the record still belongs to the same solution.
The selected record is allowed to fail the test.
It is not allowed to change the equation, switch to a neighboring solution, or replace the material packet by a detached diagnostic.
\end{definition}

\begin{lemma}[Terminal entry]
Every finite same-solution breakdown claim enters the ordered CM test before it can be used as a Clay obstruction.
\end{lemma}

\begin{proof}
A finite breakdown claim has to say that the original smooth solution reaches a terminal obstruction.
The claim therefore needs a record that can be checked against the original material motion, the original Navier--Stokes equation, and a positive-scale field readout.
Those are precisely the three tests named Pack, Part, and Field.
Without that entry, the alleged object is only a pattern associated with a proof attempt; it is not yet a terminal obstruction for the same solution.
\end{proof}

\begin{proposition}[The ordered test is exhaustive for the record]
Once a terminal record \(Q\) has entered the CM test, the record is either a member branch or the first failed face is Pack, Part, or Field.
\end{proposition}

\begin{proof}
Membership in \(\Owork\) is exactly the conjunction of the same-fluid packet, same-equation participation, and positive-scale field readout.
The test is ordered because later questions only make sense after earlier ones survive.
If Pack fails, the record has no positive material body for the original solution.
If Pack survives and Part fails, the record has a body but not the same participating equation on that body.
If Pack and Part survive and every positive Field scale fails, the record has a participating body but no coherent positive-scale field readout.
If none of these failures occurs, the record retains the faces that define membership.
There is no remaining logical slot for an in-class nonsmooth terminal branch.
\end{proof}

\section{What Each Face Means}

Pack is the material-address question.
The proof is not asking whether one can draw a small ball around a bad point.
It asks whether the alleged terminal object still contains a positive same-fluid packet whose ancestry belongs to the original solution.
A zero-radius remnant, a pure frequency shadow, or a pattern that has lost its material body fails this face.

Part is the equation-participation question.
After a packet is present, the same pressure, viscosity, and nonlinear source ancestry have to act on that packet.
This prevents the proof from proving a statement about a filtered surrogate and then silently treating it as a statement about the original Navier--Stokes evolution.

Field is the positive-scale readout question.
After the packet and the equation have survived, the record still needs one positive scale on which the field can be read coherently through the required depth.
The theorem does not require every scale to be good.
It requires enough positive-scale field structure for the alleged terminal branch to remain a member of the working class.

\begin{lemma}[Face failure is not a hidden smoothness claim]
When a terminal record fails Pack, Part, or Field, the proof is not claiming that a new smooth solution has been constructed.
It is claiming that the alleged terminal obstruction has lost the structure required to be an in-class obstruction.
\end{lemma}

\begin{proof}
The pass branch carries the continuation readout.
The fail branch carries a class-exit readout.
These are different roles.
A face failure does not have to produce a smooth extension; it has to prevent the failed record from serving as a same-solution nonsmooth member branch.
That is exactly what \(\Exit(Q;\Owork)=\neg\Member(Q;\Owork)\) records.
\end{proof}

\section{Continuation On The Pass Side}

The pass side is deliberately modest.
It does not ask for a new global bound.
It asks only for the standard local continuation trigger at the selected terminal packet.

\begin{lemma}[Pass gives continuation]
If the selected terminal record keeps Pack, Part, and a positive Field scale with an \(\Hs\), \(s>5/2\), readout up to \(T_*\), then the same solution continues past \(T_*\).
\end{lemma}

\begin{proof}
For \(s>5/2\), the usual Sobolev embedding gives the classical local well-posedness and continuation framework for smooth Navier--Stokes data.
The point of the Pack and Part requirements is that the \(\Hs\) readout belongs to the same solution, not to an exported surrogate.
The point of the Field requirement is that the readout is available at a positive scale on the selected record.
With those requirements present, the ordinary continuation criterion applies to the original branch.
Thus a pass-side terminal record cannot be terminal.
\end{proof}

\section{Whole-Space Exterior Source In Detail}

The whole-space problem adds one issue that the periodic problem does not have: an obstruction can try to hide in the exterior source field.
The proof separates that issue before using the CM test.

The Duhamel formula writes the solution as
\[
  u(t)=e^{t\nu\Delta}u_0
  -\int_0^t e^{(t-s)\nu\Delta}\mathbb P\nabla\cdot(u\otimes u)(s)\,ds.
\]
For rapidly decaying \(u_0\), the initial heat-flow tail does not create a terminal obstruction at spatial infinity.
For the fixed compact-core part of the nonlinear source, the heat kernel supplies far-field decay.
After those two pieces are removed, a remaining exterior obstruction has only two possible readings.
It either vanishes at the \(H^s\) continuation level, or it survives by escaping to finer and finer dyadic scales outside every fixed ball.

\begin{lemma}[Bounded exterior frequencies do not survive as the terminal obstruction]
After the initial heat tail and compact-core source have been removed, a genuine exterior survivor cannot remain at bounded dyadic frequency.
\end{lemma}

\begin{proof}
Bounded frequencies are spatially coarse.
Exterior \(L^2\) tightness removes such coarse mass outside sufficiently large balls once the harmless heat tail and fixed compact-core source have been separated.
Therefore a remaining exterior obstruction cannot be a bounded-frequency object.
If it survives, it must move to dyadic scales \(N_j\to\infty\) along radii \(R_j\to\infty\).
\end{proof}

\begin{lemma}[Dyadic escape is a face failure]
A dyadic exterior survivor is consumed by the Pack/Part/Field test.
\end{lemma}

\begin{proof}
If the dyadic survivor has no positive material address in the original flow, Pack fails.
If it has such an address but the same pressure-viscosity equation no longer participates on that packet, Part fails.
If Pack and Part both survive, the survivor is a high-frequency remnant on one same participating packet.
A positive Field scale would give a fixed coherent readout on that packet.
At a fixed positive readout scale, the Littlewood--Paley tail of the record vanishes as the dyadic index tends to infinity.
That contradicts the survivor lower bound along \(N_j\to\infty\).
Thus no positive Field scale remains, and the survivor lands in Field failure.
\end{proof}

\section{No Circularity}

The proof does not assume global smoothness in order to prove global smoothness.
The pass branch uses the standard local continuation criterion only after the selected record has retained the same-solution membership faces.
The fail branch does not use a positive estimate to delete the bad object.
It lets the bad object fail membership and records that failure as class exit.

\begin{proposition}[The exterior split does not smuggle in the conclusion]
The whole-space proof does not require the direct exterior \(H^s\) tail estimate as an assumption.
\end{proposition}

\begin{proof}
The direct tail estimate is one way to land on the pass side.
The proof also treats the branch where the exterior tail survives.
That survivor is not discarded.
It is tested as a terminal record.
The preceding dyadic-escape lemma sends the survivor to Pack, Part, or Field.
Thus the proof does not assume that all exterior \(H^s\) tails vanish; it proves that the alternative cannot be an in-class nonsmooth terminal branch.
\end{proof}

\section{Analytic Inputs Used By The Argument}

The proof uses four standard analytic inputs.
They are named here so that the CM part of the argument does not hide them inside new vocabulary.

\begin{enumerate}[label=(\roman*)]
\item Local well-posedness and continuation in \(H^s\), \(s>5/2\), for smooth divergence-free data.
\item Heat-kernel preservation of rapid decay for the initial linear part.
\item Far-field decay of the compact-core Duhamel source after a fixed spatial core has been selected.
\item Littlewood--Paley tail vanishing on a fixed positive field readout.
\end{enumerate}

The first input is the usual Sobolev continuation mechanism.
The second and third inputs say that the whole-space proof cannot manufacture a terminal obstruction from the harmless linear tail or from a fixed compact nonlinear core.
The fourth input is the reason that a positive Field scale is incompatible with a high-dyadic survivor.

\begin{lemma}[Fixed field readout has vanishing dyadic tail]
Let \(f\in H^s(\Rthree)\) on a fixed positive readout scale.
Then
\[
  \lim_{N\to\infty} N^s\|P_N f\|_{L^2}=0
\]
along the dyadic tail after passing to the usual square-summable Littlewood--Paley representation.
\end{lemma}

\begin{proof}
The \(H^s\) norm is equivalent to the square sum of dyadic pieces,
\[
  \|f\|_{H^s}^2 \simeq \sum_N N^{2s}\|P_N f\|_{L^2}^2.
\]
A square-summable sequence has terms tending to zero.
Therefore the dyadic \(H^s\) tail vanishes on any fixed record that genuinely has the \(H^s\) readout.
\end{proof}

\begin{lemma}[Heat and compact-core pieces are not terminal exterior obstructions]
For smooth rapidly decaying \(u_0\), the initial heat-flow part does not supply an exterior terminal obstruction.
After a fixed compact core of the nonlinear source has been selected, that compact-core Duhamel piece also does not supply the exterior terminal obstruction.
\end{lemma}

\begin{proof}
The heat kernel maps rapidly decaying smooth data to smooth functions with controlled far tails at positive times.
The fixed compact-core Duhamel term is obtained by integrating the heat kernel against a source supported, for this split, in a fixed spatial core.
Away from that core, the heat kernel gives far-field decay.
Thus neither term can be the remaining exterior object that distinguishes the whole-space terminal test from the periodic packet test.
\end{proof}

\section{The Full Proof Chain}

The proof can be read as one chain.
This section spells it out without relying on the labels alone.

\begin{proposition}[Finite terminal claim enters the chain]
Assume a finite terminal smoothness failure is claimed for the original whole-space Navier--Stokes solution.
Then the claim supplies a selected terminal record \(Q\) for the same solution.
\end{proposition}

\begin{proof}
A finite terminal smoothness failure is a claim about the solution generated by \(u_0\).
It is not merely a claim that some sequence, scale, or exterior pattern looks bad.
To be a counterexample to continuation, the bad object has to remain tied to the solution whose continuation is being denied.
That tie is the selected terminal record.
\end{proof}

\begin{proposition}[The chain has only pass or exit]
For the selected terminal record \(Q\), the CM chain has only these outcomes:
\[
\begin{array}{ll}
\text{Pass:} & Q\in\Owork\text{ and the same branch continues by the }H^s\text{ criterion},\\[3pt]
\text{Exit:} & Q\notin\Owork\text{ because the first failed face is Pack, Part, or Field.}
\end{array}
\]
\end{proposition}

\begin{proof}
The member branch retains the required same-solution faces.
The fail branch loses the first required face in the ordered test.
The ordered exhaustion lemma proves that no third same-record branch remains after these two outcomes.
The pass outcome gives continuation, and the fail outcome removes the alleged object from the class of legal same-solution terminal counterexamples.
\end{proof}

\begin{proposition}[Whole-space escape still enters the same chain]
The exterior-source branch in \(\Rthree\) does not create a third outcome.
\end{proposition}

\begin{proof}
The initial heat tail and fixed compact-core source are removed first.
The remaining exterior source either vanishes at the \(H^s\) continuation level or survives by dyadic escape.
Vanishing gives the pass-side periodic readout on the selected packet.
Dyadic escape is tested as a terminal record and lands in Pack, Part, or Field.
Thus the whole-space branch still enters the same pass-or-exit chain.
\end{proof}

\section{Referee Objections Addressed Inside The Proof}

\subsection*{Could the obstruction be a zero-radius object?}

A zero-radius object does not have a positive same-fluid packet.
That is Pack failure.
The proof does not need to pretend that such an object is smooth.
It only needs to show that the object is not a legal same-solution member branch.

\subsection*{Could the obstruction be a vortex-like point remnant?}

A vortex is not merely a label for large rotation.
As a terminal obstruction in this proof, it would still need a positive carrier, equation participation, and a positive-scale field readout.
A point-like remnant that loses positive radius fails Pack.
A remnant with a carrier but without same-equation participation fails Part.
A remnant with carrier and participation but no positive coherent field scale fails Field.

\subsection*{Could the whole-space tail avoid the local packet test?}

The exterior split is exactly what prevents that escape.
The harmless heat tail and compact-core source are removed.
The bounded-frequency exterior residue is removed by tightness.
The only survivor is dyadic escape.
Dyadic escape is not a new class of terminal object; it is tested by Pack, Part, and Field.

\subsection*{Does class exit merely rename blowup?}

No.
Class exit says that the alleged terminal record is not a member of the same-solution working class.
The Clay contradiction requires an in-class nonsmooth terminal branch for the original solution.
When the alleged object exits the class, it can no longer play that role.

\subsection*{Does the argument use a positive proof in disguise?}

The pass branch uses ordinary local continuation after the required readout is present.
The fail branch does not prove the bad object smooth.
It proves that the bad object fails the structure needed to be a legal same-solution terminal obstruction.
That is why the proof is contrapositive rather than a hidden global estimate.

\section{Reader-Level Proof Of The Main Contradiction}

Suppose a finite terminal Clay counterexample exists.
The counterexample is a statement about the original solution from the original datum.
Therefore the claimed terminal obstruction must be readable as a terminal record of that same solution.

The record is tested in order.
Pack asks for the positive material body.
Part asks whether the same Navier--Stokes equation still participates on that body.
Field asks for one positive-scale field readout on the participating body.

When all three answers are positive, the record is a member branch.
For \(s>5/2\), the positive field readout gives the standard continuation mechanism.
The solution continues, so the terminal counterexample is contradicted.

When one answer is negative, the first negative answer is the face failure.
The record is no longer a member of the working same-solution class.
The alleged obstruction has not become a legal in-class nonsmooth terminal branch.
It has become a class-exit witness.
That also contradicts the counterexample claim, because a Clay counterexample must be a terminal obstruction for the same solution, not a detached failure of the test.

The whole-space exterior source does not add a third case.
The linear tail and compact-core source are harmless for the exterior terminal question.
The bounded exterior frequencies disappear by tightness.
The only surviving exterior branch is dyadic escape, and dyadic escape is consumed by Pack, Part, or Field.

Every finite terminal counterexample claim therefore reaches either continuation or class exit.
Neither outcome is an in-class nonsmooth finite terminal branch.
The assumed finite terminal counterexample cannot exist.

\section{Formal Case Analysis}

The contradiction proof can also be written as a finite case analysis.
Let \(Q\) be the selected terminal record admitted by a claimed finite terminal counterexample.
The ordered predicate is
\[
  \Pack_Q,\qquad \Part_{N,Q},\qquad \Field_{N,r,Q}.
\]
The cases are not optional proof branches.
They are the only ways the record can be read after it enters the CM test.

\subsection*{Case 1: Pack fails}

The record has no positive same-fluid material packet.
The alleged terminal object is therefore not a terminal object of the same solution in the sense needed for the Clay contradiction.
It may be a diagnostic shadow of a failed construction, but it is not a same-solution member branch.
The conclusion is
\[
  \neg\Pack_Q \Longrightarrow \neg\Member(Q;\Owork).
\]
The proof stops at Pack because Part and Field are not meaningful without a packet.

\subsection*{Case 2: Pack survives and Part fails}

The record has a positive material packet, but the same Navier--Stokes equation no longer participates on that packet through the required depth.
This is the case that blocks surrogate drift.
A filtered object, an exported pressure shadow, or a neighboring equation may carry information, but it cannot be used as the same solution unless the pressure, viscosity, and nonlinear source ancestry still participate on the packet.
The conclusion is
\[
  \Pack_Q\wedge\neg\Part_{N,Q}
  \Longrightarrow \neg\Member(Q;\Owork).
\]
The proof stops at Part because a field readout from the wrong equation is not a field readout of the original terminal branch.

\subsection*{Case 3: Pack and Part survive and Field fails}

The record has a positive same-fluid packet and the same equation participates on it.
The remaining question is whether any positive scale carries a coherent field readout.
When every positive scale fails,
\[
  \forall r>0\,\neg\Field_{N,r,Q},
\]
the record still cannot be a member branch.
The terminal obstruction has no positive-scale field window on which the continuation-level readout can be made.
The conclusion is
\[
  \Pack_Q\wedge\Part_{N,Q}\wedge
  \forall r>0\,\neg\Field_{N,r,Q}
  \Longrightarrow \neg\Member(Q;\Owork).
\]

\subsection*{Case 4: Pack, Part, and Field survive}

The record remains a same-solution member branch.
The Field readout is at a positive scale and at Sobolev level \(s>5/2\).
The standard continuation theorem applies to the original branch.
Thus
\[
  \Pack_Q\wedge\Part_{N,Q}\wedge\exists r>0\,\Field_{N,r,Q}
  \Longrightarrow \Member(Q;\Owork)
  \Longrightarrow \text{same-solution continuation}.
\]

These four cases exhaust the record.
The first three are class exit.
The fourth is continuation.
No case gives an in-class nonsmooth terminal branch.

\section{Interface With The Classical Continuation Criterion}

The proof does not replace local well-posedness.
It uses local well-posedness at the one place where that theorem belongs.
For \(s>5/2\), the classical theory says that an \(H^s\) solution can be continued as long as the \(H^s\) control needed for the local theory remains available.
The CM proof supplies that control only on the pass branch.

This matters because a direct proof would try to estimate every dangerous nonlinear term until the global \(H^s\) norm is bounded.
The present proof does not take that route.
It says that a terminal counterexample has to present a terminal record.
When that record retains the \(H^s\) continuation readout, the classical theory continues it.
When the record does not retain that readout, the reason for failure is classified by Pack, Part, or Field.

\begin{lemma}[The continuation criterion is used only after same-solution custody]
The \(H^s\) continuation criterion is invoked only for records that have retained Pack and Part.
\end{lemma}

\begin{proof}
Pack gives the positive material carrier for the original solution.
Part gives the participation of the original Navier--Stokes equation on that carrier.
Without these two faces, an \(H^s\) statement might belong to the wrong object.
The proof therefore invokes the continuation criterion only after the readout has been tied to the same solution.
\end{proof}

\section{Whole-Space Localization Without Changing The Problem}

The whole-space argument must not become a periodic theorem in disguise.
The selected packet is local, but the equation remains the whole-space equation.
The pressure projection, the heat kernel, and the nonlinear source are still the \(\Rthree\) objects.
The localization is used only to ask whether the alleged terminal obstruction remains attached to a positive same-fluid packet.

The exterior split is therefore a custody device, not a theorem change.
The initial heat tail is separated because rapid decay makes it harmless for terminal obstruction.
The compact-core source is separated because heat-kernel propagation from a fixed core cannot create the remaining exterior terminal survivor.
The exterior source is then tested honestly.
It either vanishes at the continuation level or survives by dyadic escape.

\begin{proposition}[Localization preserves the target equation]
The Pack/Part/Field test on a whole-space terminal packet does not replace the \(\Rthree\) Navier--Stokes equation by a periodic equation.
\end{proposition}

\begin{proof}
The periodic proof supplies the local terminal-witness grammar.
The whole-space proof uses that grammar only after the whole-space exterior terms have been separated.
Part requires the same whole-space pressure-viscosity participation on the selected packet.
Therefore the packet test cannot silently become a proof about a different equation.
\end{proof}

\section{The Dyadic Survivor Cannot Be A Member Branch}

The surviving exterior source is the only branch that needs special attention.
It is not enough to say that the survivor is far away.
A far-away object could still matter if it remained attached to the same solution.
The proof instead asks what structure the survivor has.

A bounded-frequency survivor has spatial width.
After the harmless pieces have been removed, exterior tightness eliminates such a survivor outside large balls.
The remaining survivor must therefore be high dyadic frequency.
High dyadic frequency gives the branch its pressure: the object can keep showing up along \(N_j\to\infty\).
The same fact also prevents it from being a fixed positive Field readout.

\begin{proposition}[Dyadic escape contradicts fixed positive Field]
A dyadic exterior survivor with \(N_j\to\infty\) cannot also be the same fixed positive-scale \(\Field_{N,r,Q}\) readout of a member record.
\end{proposition}

\begin{proof}
Assume Pack and Part survive, so the survivor is being tested on one same participating packet.
Assume also that a fixed positive Field readout survives.
Then the field readout is an \(H^s\) object on that fixed record.
The dyadic tail of an \(H^s\) object tends to zero.
A survivor lower bound along \(N_j\to\infty\) says that the dyadic tail does not vanish along the selected exterior sequence.
The two statements cannot both hold for the same fixed record.
Therefore Field fails.
\end{proof}

\section{What The Main Theorem Proves}

The theorem proves that a finite terminal smoothness failure cannot remain as a legal same-solution counterexample.
That is the exact Clay target in this paper.
The theorem does not say that every bad-looking diagnostic is smooth.
It says that every alleged terminal obstruction either gives the continuation readout or loses the structure required to be the original solution's terminal obstruction.

This is why the class-exit conclusion is proof-bearing.
The Clay counterexample needs an object that is both terminal and still the original solution's object.
The pass branch destroys terminality by continuation.
The fail branch destroys same-solution membership by the first failed face.
The assumed counterexample has no branch left.

\section{Expanded Proof Of The Periodic Chain}

The periodic chain is the clean model because there is no exterior infinity.
The alleged obstruction is forced to live on the torus with the same fluid.
The proof starts with a terminal record \(Q\) selected from the original smooth branch.
The selection is not a new solution.
It is the information by which the alleged finite terminal failure is supposed to be recognized.

The first question is Pack.
The torus has no spatial infinity where the obstruction can hide, but a terminal record can still fail to provide a positive material packet.
For example, a sequence of diagnostics can shrink until no positive same-fluid body remains.
Such a sequence may be useful for describing why a direct estimate failed.
It is not a same-solution member branch.
The proof records this as Pack failure.

When Pack survives, the record has a positive body.
The next question is Part.
The body must be acted on by the same Navier--Stokes equation.
The pressure term matters here.
The nonlinear source and pressure are nonlocal, so it is not enough to keep a visual packet.
The packet has to keep the same pressure-viscosity participation that the original equation supplies.
When this participation fails, the record cannot be the original solution's terminal branch.

When Pack and Part survive, the record has a participating body.
The final question is Field.
The alleged terminal failure now has to be visible on a positive scale.
The proof does not demand that every scale be controlled.
It demands that the record retain at least one positive field window that can carry the continuation-level readout.
When every positive field window fails, the alleged object has no coherent terminal readout inside the working class.

The periodic theorem follows from these three questions.
The same terminal record either passes the test or fails at its first failed face.
On the pass side, the record carries the same-solution \(H^s\) continuation readout.
On the fail side, the record is excluded from the working class.
The same record cannot both retain all membership faces and fail one of them.
Thus the periodic problem has no legal in-class nonsmooth terminal branch.

\section{Expanded Proof Of The Whole-Space Chain}

The whole-space chain starts with the periodic grammar and adds the exterior source split.
The selected terminal packet is local, but the theorem is still the \(\Rthree\) theorem.
That is why the Duhamel decomposition is used before the CM test is applied to the exterior issue.

The linear heat part is controlled by rapid decay of the datum.
It does not create a new nonlinear terminal obstruction at infinity.
The compact-core Duhamel part is controlled by the heat kernel once the core is fixed.
It also does not create the remaining exterior terminal obstruction.
After those two pieces are removed, the only exterior issue is the nonlinear source that stays genuinely exterior.

There are two branches.
On the first branch, the exterior \(H^s\) tail vanishes.
Then the whole-space record reads like the periodic packet at the selected terminal scale.
The same Pack/Part/Field test applies, and the pass branch gives continuation.

On the second branch, the exterior source survives.
The survivor cannot be bounded-frequency exterior mass, because bounded exterior frequency is removed by tightness after the harmless pieces have been separated.
The survivor therefore escapes to dyadic scale.
This dyadic escape still has to be the same solution's terminal obstruction in order to be a Clay counterexample.
So the proof tests it by Pack, Part, and Field.

The dyadic survivor fails Pack when it has no positive material address.
It fails Part when the material address exists but the same equation does not participate on it.
It fails Field when Pack and Part survive but no fixed positive field window can contain the high-dyadic survivor.
The Field case is the decisive high-frequency case: a fixed \(H^s\) readout has dyadic tail tending to zero, while the survivor branch asserts a nonvanishing dyadic lower bound.

Thus the whole-space exterior branch is also exhausted.
Vanishing exterior tail returns to the pass-side continuation readout.
Surviving exterior tail exits by Pack, Part, or Field.
No whole-space branch produces a same-solution terminal object that remains in the class while failing smooth continuation.

\section{Why The Argument Is A Contrapositive Proof}

The direct proof style would try to show that no dangerous term can grow.
This paper does something different.
It lets the dangerous terminal object appear and asks what it is allowed to be.
The contradiction comes from the object's required identity.

The object must be terminal.
It must be terminal for the original solution.
It must still have enough structure to be a same-solution branch.
It must also fail continuation.
The Pack/Part/Field test shows that these requirements cannot all hold.

The pass branch keeps the same-solution structure and therefore continues.
The fail branch loses the same-solution structure and therefore is not a legal terminal counterexample.
The proof does not need to smooth the failed object.
It only needs to prevent the failed object from being the original solution's in-class terminal branch.

This is the point at which the proof differs from a failed positive estimate.
A failed estimate often leaves the reader with a bad scenario and no way to classify it.
Here the bad scenario is the object under test.
The worse it behaves, the more sharply it lands in Pack, Part, or Field.

\section{Submission-Scale Dependency List}

The main theorem depends on the following statements, all of which are stated in the paper.

\begin{enumerate}[label=\arabic*.]
\item A finite Clay counterexample has to supply a same-solution terminal record.
\item A same-solution terminal record is tested by Pack, Part, and Field.
\item The first failed face gives \(\Exit(Q;\Owork)=\neg\Member(Q;\Owork)\).
\item A pass-side record carries the \(H^s\), \(s>5/2\), continuation readout.
\item The whole-space exterior source either vanishes at the continuation level or survives by dyadic escape.
\item A dyadic exterior survivor fails Pack, Part, or Field.
\item Therefore no terminal record is both in-class and nonsmooth terminal.
\end{enumerate}

This list is included to make the paper checkable.
A reader can reject the proof only by rejecting one of these dependencies.
The proof does not ask the reader to trust an unpublished branch ledger or a hidden computation.
The long companion material explains how these dependencies were discovered and pressure-tested.
The submitted paper is the chain above.

\section{Referee Verification Walkthrough}

This section follows the proof as a referee check rather than as a discovery narrative.
Start with the negation of the theorem.
There is smooth rapidly decaying divergence-free data \(u_0\), a fixed viscosity \(\nu>0\), and a finite time \(T_*\) at which the same smooth solution cannot be continued.
The proof asks the counterexample to identify its terminal obstruction.

The first verification point is identity.
The object under discussion has to belong to the original solution.
An object that is merely close to the solution, derived from a filtered equation, or obtained by changing the pressure relation is not enough.
That is why Pack and Part appear before Field.
They prevent the proof from using a good estimate on the wrong object or a bad pattern that is no longer the original solution's pattern.

The second verification point is scale.
A terminal obstruction has to be visible at some positive scale if it is going to be a terminal branch of the same solution.
Without a positive scale, the obstruction has no field window on which the continuation failure can be read.
That is why a zero-radius remnant is not treated as a mysterious third branch.
It is a failure of the Field question after the earlier faces have been checked, or a Pack failure when no material body remains.

The third verification point is continuation.
When the field window exists at \(H^s\), \(s>5/2\), the proof does not invent a new continuation theorem.
It calls the standard continuation theorem.
The only work done by the CM proof is to make sure that the readout belongs to the same solution.

The fourth verification point is the whole-space exterior.
The proof does not ignore infinity.
It separates the linear heat tail, the compact-core nonlinear source, and the genuinely exterior nonlinear source.
The first two pieces cannot be the remaining terminal exterior obstruction.
The third piece is either harmless at the continuation level or survives by dyadic escape.
That survivor is then tested by the same three faces.

The fifth verification point is contradiction.
The pass branch contradicts terminality because the solution continues.
The fail branch contradicts counterexample legality because the alleged obstruction has exited the same-solution class.
The assumed finite terminal Clay counterexample needs both terminality and legality.
The proof removes one of them in every case.

\section{Notation And Logical Roles}

The notation in the paper has only a few logical roles.
The roles are separated here to prevent a common misreading.

\begin{center}
\begin{tabular}{L{0.25\linewidth}L{0.65\linewidth}}
Symbol & Role in the proof\\
\hline
\(Q\) & Selected terminal record extracted from the original solution claim.\\
\(\Owork\) & Working same-solution class: packet, participation, and field readout.\\
\(\Pack_Q\) & Positive same-fluid material carrier for \(Q\).\\
\(\Part_{N,Q}\) & Same Navier--Stokes participation on the selected carrier through depth \(N\).\\
\(\Field_{N,r,Q}\) & Positive-scale coherent field readout after Pack and Part have survived.\\
\(\Member(Q;\Owork)\) & The selected terminal record remains in the working same-solution class.\\
\(\Exit(Q;\Owork)\) & The selected terminal record exits the class, meaning \(\neg\Member(Q;\Owork)\).\\
\end{tabular}
\end{center}

Pack, Part, and Field are not synonyms for class exit.
They are faces of the test.
Class exit is the conclusion supported by the first failed face.
Likewise, \(\Member(Q;\Owork)\) is not a generic statement that a smooth solution exists.
It is the statement that the selected terminal record keeps the same-solution faces required for the continuation readout.

\section{Final Logical Compression}

The proof can be compressed to the following implication chain:
\[
\begin{aligned}
\text{finite Clay terminal counterexample}
&\Longrightarrow \text{same-solution terminal record }Q\\
&\Longrightarrow
\left[
\begin{array}{l}
\Pack_Q\wedge\Part_{N,Q}\wedge\exists r>0\,\Field_{N,r,Q},\\
\text{or the first failed face is Pack, Part, or Field}
\end{array}
\right]\\
&\Longrightarrow
\left[
\begin{array}{l}
\text{same-solution continuation},\\
\text{or }\Exit(Q;\Owork)
\end{array}
\right].
\end{aligned}
\]

The first outcome contradicts finite terminal breakdown.
The second outcome contradicts the existence of a legal in-class terminal counterexample.
Those are the only outcomes of the selected record.
Therefore the original finite terminal counterexample assumption is false.

\section{What Is Excluded}

The proof excludes a specific object: an in-class finite terminal nonsmooth branch of the original whole-space Navier--Stokes solution.
That is the object required by a finite-time Clay counterexample.
The proof does not need to classify every diagnostic pattern that can appear in a failed estimate.
It only needs to show that none of those patterns can remain as the same solution's legal terminal obstruction.

A bad object with no positive material carrier is excluded at Pack.
A bad object with a carrier but without the same pressure-viscosity participation is excluded at Part.
A bad object with carrier and participation but no positive field window is excluded at Field.
A record with carrier, participation, and a positive \(H^s\) field window continues.
The alternatives are exhaustive because the working class is exactly built from those faces.

This is also why the proof can tolerate violent-looking terminal scenarios.
The worse scenario does not have to be smoothed by force.
It has to identify how it remains the original solution's terminal branch.
When it cannot do that, it exits the class.
When it can do that and carries the continuation readout, it is not terminal.

\section{What The Proof Does Not Rely On}

The proof does not rely on a global a priori bound for every derivative.
It does not rely on proving that the exterior nonlinear source always vanishes.
It does not rely on replacing the whole-space problem by a periodic problem.
It does not rely on treating a support diagnostic as the theorem.

The proof relies instead on terminal object custody.
The counterexample must supply a same-solution terminal record.
That record is either a continuation record or a class-exit record.
The exterior source branch is forced into the same decision by the Duhamel split and the dyadic survivor test.

The positive information used on the pass side is exactly the local continuation readout.
The negative information used on the fail side is exactly the first failed membership face.
These two uses do not collapse into each other.
Continuation is a statement about the member branch.
Exit is a statement about the failed record.

\section{Submission Form Of The Result}

The submitted theorem can be stated without the internal history of how the proof was found.

\begin{quote}
Every smooth rapidly decaying divergence-free whole-space Navier--Stokes datum has no finite same-solution terminal nonsmooth branch.
Any alleged terminal branch either retains the \(H^s\), \(s>5/2\), continuation readout or fails the same-solution class-membership test at Pack, Part, or Field.
Therefore finite-time loss of smoothness cannot occur.
\end{quote}

The rest of the paper is the proof of this statement.
The early sections define the terminal record and the working class.
The middle sections prove that the ordered faces exhaust the terminal record.
The whole-space sections prove that exterior escape still lands in the same test.
The final sections convert pass-or-exit into the Clay contradiction.

\section{Final Theorem Check}

At the end of the proof there is only one possible place for a counterexample to hide.
It would have to be a terminal record that is still the original solution's record, still acted on by the original equation, still visible on a positive field scale, and still unable to continue.
The paper has already treated that exact object.

The first three requirements are Pack, Part, and Field.
When the record keeps them, the \(H^s\) continuation readout applies.
When the record loses one of them, the record is outside the same-solution working class.
The counterexample needs both membership and terminal nonsmoothness.
The proof makes those two requirements incompatible.

The whole-space issue does not reopen the hiding place.
The exterior source is separated before the final test.
A vanished exterior tail leaves the local packet continuation question.
A surviving exterior tail is dyadic escape.
Dyadic escape cannot be a fixed positive Field readout on a same participating packet.
It therefore joins the same face-failure alternatives.

Thus the final theorem has exactly the advertised scope.
It is not a periodic-only result.
It is not a status statement about a proof program.
It is the whole-space finite-terminal counterexample exclusion for smooth rapidly decaying data.

\section{From Class Exit To The Clay Statement}

\begin{definition}[Clay terminal counterexample]
Fix \(s>5/2\).
A Clay terminal counterexample in this paper means a smooth rapidly decaying divergence-free datum \(u_0\) and a finite time \(T_*\) for which the same Navier--Stokes solution from \(u_0\) cannot be continued as an \(\Hs\) solution past \(T_*\).
The obstruction is not allowed to be a detached pattern at infinity, a filtered surrogate, or a neighboring equation.
It must be a terminal record of the same solution.
\end{definition}

\begin{lemma}[A Clay obstruction needs an in-class terminal record]
Any finite-time Clay terminal counterexample supplies a terminal record \(Q\) that remains eligible for the same-solution class-membership test.
In particular, a detached object that has already lost Pack, Part, or every positive Field scale is not a legal in-class nonsmooth counterexample.
\end{lemma}

\begin{proof}
The Clay question asks about the smooth solution generated by one datum and one Navier--Stokes equation.
A terminal obstruction must therefore still be about that solution.
It has to keep a positive material packet, the participation of the same pressure-viscosity equation on that packet, and a positive-scale field readout on which the continuation failure is being claimed.
Those are exactly the Pack, Part, and Field faces of the class-membership test.
When one of those faces has already failed, the object may record where the alleged obstruction broke, but it no longer gives a same-solution member branch that can be both terminal and nonsmooth inside the working class.
\end{proof}

\begin{theorem}[Clay smoothness conclusion]
For every smooth rapidly decaying divergence-free datum on \(\Rthree\) and every fixed \(\nu>0\), the three-dimensional incompressible Navier--Stokes solution does not terminate at a finite time by loss of smoothness.
\end{theorem}

\begin{proof}
Assume for contradiction that a smooth rapidly decaying datum produces a finite terminal Clay counterexample at time \(T_*\).
By the preceding lemma, the claimed obstruction supplies a same-solution terminal record \(Q\) eligible for the CM test.

The whole-space CM completion theorem gives only two outcomes.
On the pass side, the selected packet keeps the continuation-level \(\Hs\) readout with \(s>5/2\), so the standard local continuation criterion extends the same solution past \(T_*\).
That contradicts terminality.

On the fail side, the first failed face is Pack, Part, or Field, and the CM embedding gives
\[
  \Exit(Q;\Owork)=\neg\Member(Q;\Owork).
\]
Then \(Q\) is not an in-class same-solution nonsmooth terminal branch.
That contradicts the requirement that the alleged Clay obstruction be a legal terminal counterexample for the original solution.

Both branches contradict the assumed finite terminal counterexample.
Therefore no such finite terminal loss of smoothness exists.
\end{proof}

\section{Why The Long Dossier Is Not The Submission Paper}

The companion work surface also contains a long app-aligned dossier.
That dossier records branch history, failure analysis, source-field expansion, and graph-facing proof roles.
It is useful for inspection and for rebuilding the proof surface.
It is not the shape of this paper.

This manuscript is the Clay-facing argument.
The proof burden is concentrated in the same-solution terminal test:
the terminal record enters the CM test, the ordered Pack/Part/Field faces consume the fail branch, the pass branch retains the \(H^s\) continuation readout, and the whole-space exterior survivor is either harmless tail or a face failure.
The long material remains provenance and working support outside the submitted paper.

\section{Reader Check}

A referee should be able to test the proof by asking four questions.
First, is the alleged obstruction still a same-solution terminal object?
Second, does it retain a positive material packet?
Third, does the same Navier--Stokes equation still participate on that packet?
Fourth, is there a positive field scale on which the terminal readout is coherent?

The pass branch answers all four questions positively and gives continuation.
The fail branch answers one of the ordered questions negatively and leaves the working class.
The whole-space argument adds only the exterior-source split: the exterior source either vanishes at the continuation level or survives only as dyadic escape, and dyadic escape has to fail Pack, Part, or Field.

\section{Provenance Boundary}

The separate app-aligned dossier is the expanded work record.
This Codex paper uses that work record as source material, but it does not import the dossier pages into the submission PDF.
The theorem claims above are meant to stand on the page as the submitted argument.

\end{document}
